\section{Related Work}
\label{sec:related}

We organize prior work by defense stage and highlight the
structural gap that motivates Security Cuff.

\paragraph{Visual backdoor attacks on VLA.}
VLA backdoor attacks share a common mechanism: an adversary poisons
training demonstrations to bind a visual trigger to a malicious
action sequence.
They differ in trigger design and stealth.
Physical-object triggers~\citep{goba} activate when an everyday
item enters the scene; visual-patch triggers~\citep{tabvla} are
injected via black-box fine-tuning; optimized pixel
perturbations~\citep{badvla} are imperceptible to human inspection.
AttackVLA~\citep{attackvla} provides a systematic benchmark
covering both adversarial and backdoor vectors with bi-modal
triggers.
At the stealth frontier,
SilentDrift~\citep{silentdrift} crafts $C^2$-continuous trajectory
deviations that exploit action chunking, and
State Backdoor~\citep{statebackdoor} eliminates visual artifacts
entirely by using the robot's initial state as trigger.
Similar vulnerabilities appear in VLM-based
agents~\citep{beat,trojanrobot} and LLM-based
agents~\citep{contextbackdoor}.
These attacks preserve clean behavior while steering a subset of
triggered rollouts toward harmful trajectories, which is why our
defense focuses on runtime consequence verification rather than
trigger reconstruction.

\paragraph{Pre-deployment defenses.}
Training-time methods---spectral
signatures~\citep{tran2018spectral}, activation
clustering~\citep{chen2019detecting},
fine-pruning~\citep{liu2018fine}---detect or remove backdoors but
require centralized access to the full training set, an assumption
that breaks down as VLA systems adopt cloud-edge continuous learning
pipelines where weights are periodically fine-tuned on
edge-collected data.
Input-side screening operates at test time:
STRIP~\citep{gao2019strip} measures perturbation sensitivity;
SampDetox~\citep{sampdetox} purifies inputs via diffusion-based
denoising.
Model-inspection tools such as Neural
Cleanse~\citep{wang2019neural} reverse-engineer triggers from a
trained model.
SafeVLA~\citep{safevla} incorporates safety constraints during
training but does not target deliberately planted backdoors whose
triggers are absent from the alignment data.
All of these methods operate \emph{upstream} of the model's output
actions, whereas Security Cuff intervenes after policy execution has
already begun.

\paragraph{Inference-time safety monitors.}
Most runtime approaches target \emph{accidental} failures rather
than deliberate backdoor attacks.
Output-uncertainty methods~\citep{mahal,ensemble_ood} flag high
predictive entropy or OOD inputs; a capable adversary can suppress
action-level anomalies~\citep{silentdrift}, making such signals
unreliable as a sole defense.
SAFE~\citep{safe} trains probes on VLA hidden states to predict
task failure; however, a backdoored model executes the attack with
internal coherence resembling successful task execution, so the
failure probe detects no anomaly.
Constraint-based filters such as VLSA~\citep{vlsa} enforce
geometric constraints via CBF~\citep{ames2017cbf}, but cannot
intercept actions that are geometrically safe yet semantically
dangerous---the dominant pattern in VLA backdoor attacks.
BERA~\citep{bera} reconstructs visual tokens to erase visual
triggers before policy execution; however, this pre-encoding
intervention cannot detect triggers that manifest through hidden-state
representations rather than localizable visual
artifacts~\citep{statebackdoor,silentdrift}.
Together, these gaps motivate a post-encoding,
consequence-aware defense that reads internal model signals and
verifies predicted physical outcomes rather than searching only for
visual artifacts or raw uncertainty.

\paragraph{World models in embodied robotics.}
Learned world models~\citep{ha2018world,daydreamer} have enabled
imagination-based planning and latent rollouts in robotics.
Recent world action models (WAMs), such as
World Action Models are Zero-shot Policies~\citep{wamzero} and
Self-Correcting VLA~\citep{scvla}, push this idea closer to control by
using imagined futures to choose, chunk, or refine actions during
execution.
Security Cuff uses the same predictive machinery differently: it does
not use rollout to improve the next action or the policy itself.
Instead, the base policy is left fixed and the world-model rollout is
used only as an external consequence-aware verifier that decides
whether a proposed prefix is likely to lead to catastrophic physical
outcomes.
